\section{统计学引论}

这门课的第一章不讲概率，先讲一件事：数据是怎么变成结论的，以及中间有哪些环节会把结论弄坏。看起来像废话，但后面每一章的假设检验、置信区间，判定的都是第一章这几个毛病。

\subsection{统计学是什么}

\begin{definition}[统计学]
统计学是从数据中学习的技艺，关心三件事：
\begin{enumerate}
 \item 数据的收集；
 \item 数据的描述与整理；
 \item 数据的分析——并由此得出结论。
\end{enumerate}
\end{definition}

\begin{proposition}[四步是一条单向流水线]
收集数据 $\to$ 描述数据 $\to$ 分析数据 $\to$ 得出结论。每一步都喂给下一步，最后一步必须是「站得住的结论」：能被别人用同一批数据复查的那种。
\end{proposition}

链条上任何一环糊弄了，后面再漂亮也救不回来。数据收集阶段埋的坑，到分析阶段是看不出来的——这大概是整门课最想让人记住的一句话。

\subsection{数据从哪来}

统计工作只有两种起点。

\begin{enumerate}
 \item \textbf{数据已经存在}：政府每年公布降水、地震次数、失业率、GDP、通胀。要做的只是描述、整理、分析。
 \item \textbf{数据还不存在}：得自己设计实验把它们造出来。
\end{enumerate}

\begin{example}[两种教学法]
想知道哪种编程教学法更好，把学生分成两组，各用一种方法教，最后比考试成绩。
\end{example}

这是最典型的第二种起点，而它成败的地方在分组：

\begin{proposition}[分组必须随机分配]
不能用「男生一组、女生一组」，也不能用任何可能与能力相关的规则分组。正确做法是随机分配，使得每一种可能的分组方式出现的概率相同。
\end{proposition}

理由不是「随机比较高级」，而是只有随机分配才能把「两组本来就不一样」这件事变成可以用概率算的东西。按成绩好坏自己挑着分，差异里就掺进了选择规则本身，后面怎么算都说不清。

\begin{remark}[「等概率」是哪种设计的要求]
「每一种可能的分组方式等概率」是\emph{完全随机分组}的要求，不是所有随机实验设计的通用定义。随机区组、分层设计、拉丁方各自排除了一部分分组方式，被排除的那些概率是零。随机化真正要求的是：分组规则由概率机制决定、不掺研究者的偏好，且每个单元入组的概率事先已知。
\end{remark}

\begin{definition}[描述统计]
描述统计负责数据的描述与汇总，例如把两组分数列出来，或者各算一个平均分。
\end{definition}

\subsection{推断统计与概率模型}

\begin{definition}[推断统计]
推断统计负责从数据中得出结论。
\end{definition}

麻烦在于：看到差异，不能直接归因。

\begin{proposition}[差异可能只是偶然]
第一组平均分高出一截，可能是教学法真有用，也可能纯属运气。推断统计要做的，就是把「真差异」和「偶然波动」区分开。
\end{proposition}

\begin{example}[硬币]
一枚硬币抛 10 次出现 7 次正面，很可能只是一枚普通硬币碰巧如此；抛 50 次出现 47 次正面，直觉上就很难再用「碰巧」搪塞了。同样是「正面偏多」，样本量不同，说服力完全不同。
\end{example}

\begin{remark}[直觉结论和推断结论不是一回事]
「47 次正面所以硬币不均匀」只是直觉。要写成正式结论，得先把原假设（$p=1/2$）、备择假设（双侧还是单侧）、显著性水平，以及「各次独立、每次正面概率相同」这些模型假设摊开，才算一个可复查的检验。脱离这几样，「几乎可以断定」这句话没有确定含义。
\end{remark}

要把这件事说清楚，就得提前对「各种结果出现的可能性」做假设。

\begin{definition}[概率模型]
为了从数据得出结论而对数据各取值的可能性所作的全部假设，称为这批数据的概率模型。
\end{definition}

注意定义里的位置：概率模型是\emph{手段}，是为了推断才搭的脚手架，不是数据本身的属性。模型假设错了，后面的结论就只是逻辑自洽而已。

\begin{example}[广告点击率]
想知道新广告的用户点击比例 $p$，就把广告随机展示给一批用户并记录点击。假设每个用户点击的概率都是 $p$，$p$ 未知，用观察到的数据去推断它。
\end{example}

\subsection{总体与样本}

\begin{definition}[总体与样本]
总体是想要了解的全部对象的集合；样本是其中真正被观测的那个子集。
\end{definition}

例如「某州全部居民」「去年出厂的全部电视机」「某社区的全部家庭户」都是总体。总体通常太大，没法逐个查，于是抽样本、观测样本，再反过来推断总体。整套流程就一张图：

\[
\text{总体}\ \xrightarrow{\ \text{随机抽样}\ }\ \text{样本}\ \xrightarrow{\ \text{统计推断}\ }\ \text{关于总体的结论}.
\]

样本要能代表总体，结论才有意义。

\begin{example}[商场门口问收入]
想估计全市家庭的平均月收入，于是在工作日下午，在一家高档商场门口拦访前 100 位顾客。这个平均值大概率偏——会在那种场合、那个时间段闲逛的人，收入结构和作息都跟全市人口不一样。
\end{example}

\begin{proposition}[随机的悖论]
完全随机地抽，是避免人为选择偏差最省事的办法。它的好处不在于「偶然比人聪明」，而在于：随机抽样给出每个单元事先已知的入样概率，于是抽样误差本身是可以算的。
\end{proposition}

\begin{remark}[分层随机抽样不是配额抽样]
简单随机抽样不是唯一可靠的概率抽样设计。总体有明显层次（性别、年龄段、职业、地区）时，常先按这些变量把总体分成不重叠的层，再在每层内部随机抽，这就是分层随机抽样。它一边保住随机性，一边保证每层都抽到足够样本；在「层内比层间均匀得多」的常见情形下，它的估计方差还往往比简单随机抽样小一截。

真正要避免的是配额抽样：让访问员「男的女的各找 50 个」，比例是凑对了，可每一格里挑谁没有概率机制，入样概率无从谈起。配额抽样和分层随机抽样表面像，差的正是层内那一遍随机。

所以不能把「人为控制样本构成」一概当成错误。该提防的是没有随机机制的人为凑数，而不是控制本身。
\end{remark}

样本抽好之后，才轮到统计推断出场。

\subsection{一点点历史}

「statistics」这个词来自 state，最早指对国家有用的事实汇编。十七世纪公共卫生很差，死亡记录是刚需：

\begin{itemize}
 \item 1562 年起伦敦每周公布死亡清册（bills of mortality），最初是瘟疫预警，好让有钱人及时出城；
 \item Graunt（1662）从清册里算出「大约每 88 人死 1 人」，用 13{,}200 的死亡数推出伦敦人口约 38.7 万，还按年龄整理了死亡表；
 \item Halley（1693，发现彗星那位）用死亡表算出各年龄的人活到各年龄的概率，说服保险人：年金保费应该按年龄定；
 \item 十九世纪的 Galton 与 Pearson 做出回归、相关与 $\chi^2$ 拟合优度检验，Gosset（笔名 Student）做出 $t$ 检验，Fisher 与 Neyman、Pearson 建起现代推断的框架。
\end{itemize}

在此之前，概率论虽然已经被 Bernoulli、Gauss、Laplace 发展起来了，统计学家却几乎不用它——他们宁愿让数据自己说话，宁可普查也不抽样。这两条线直到十九世纪末才合流。这几位后面都会重逢：$t$ 检验、$\chi^2$ 检验、回归各占一章。

\begin{example}[Graunt 的算法]
1658 年伦敦清册记录死亡 12{,}246 人，教区调查显示当年死亡率约 $2\%$。用观测死亡数除以死亡率：
\[
\text{人口}\approx\frac{12{,}246}{0.02}=612{,}300 .
\]
这就是「用一个已知的比率，把看不见的总量反推出来」，和后来所有估计问题的骨架是同一个。
\end{example}

\begin{remark}[这一节的史料出处]
Graunt 的观察出自《关于死亡清册的自然与政治观察》（1662），Halley 的年金表登在 Royal Society 的 \emph{Philosophical Transactions}（1693）。上面的数字是二手转述，「伦敦」的范围、清册漏报多少，各本统计史的口径并不一致，别当精确史料引用；要引就引原始文献，别引这份笔记。
\end{remark}

\subsection*{本章容易踩的地方}

\begin{itemize}
 \item 把「随机抽样」理解成「随便挑」。随机有明确的等可能性要求，随手抓一批往往带着挑人的规则。
 \item 只认简单随机抽样，把分层随机抽样一并当成「人为凑比例」排除掉。该排除的是没有概率机制的配额抽样。
 \item 忽略无应答：抽了 200 人只回来 86 份，这 86 份的平均值不能直接当总体平均——没回的人可能系统性地不一样。
 \item 看见组间差异就归因到处理上，忘了偶然本身也能造出差异。样本量决定同一现象值不值得惊讶。
 \item 把直觉上的「肯定有问题」当成检验结论。不说清假设与显著性水平，这类断言无法复查。
\end{itemize}
